% ======================================================================================
% Nom     : reports/latex/sections/09_limites_ameliorations.tex
% Rôle    : Section présentant les limites et les améliorations possibles.
% Auteur  : Maxime BRONNY
% Version : V1
% Cadre   : Réalisé dans le cadre du cours Techniques d'apprentissage artificiel
%           Master 1 Informatique Big Data
% Usage   : Fichier inclus dans main.tex lors de la compilation du rapport.
% ======================================================================================

\section{Limites et améliorations possibles}
\label{sec:limites}

Les limites identifiées dans l'analyse critique dessinent directement les
améliorations envisageables, classées ici de la plus simple à la plus
structurante.

\paragraph{Encodage one-hot des variables nominales.}
Remplacer le label encoding de \texttt{loan\_intent} et
\texttt{person\_home\_ownership} par un encodage one-hot (une colonne binaire
par modalité) supprimerait l'ordre artificiel entre modalités. Le bénéfice
attendu concerne d'abord la régression logistique, qui pourrait attribuer un
poids propre à chaque modalité ; le coût est minime (quelques colonnes
supplémentaires).

\paragraph{Régularisation L2.}
Notre régression logistique ne comporte aucune pénalité sur les poids, là où
scikit-learn en applique une par défaut. C'est l'une des sources du petit
écart observé. Ajouter un terme $\lambda \lVert \mathbf{w} \rVert^2$ au coût
(et $2\lambda\mathbf{w}$ au gradient) est immédiat et stabiliserait le modèle,
notamment après un passage au one-hot qui augmente le nombre de colonnes.

\paragraph{Validation croisée $k$-fold.}
Le découpage unique $70/15/15$, même stratifié et à graine fixe, fournit une
estimation ponctuelle des performances. Une validation croisée à $k$ blocs
donnerait une moyenne et un écart-type pour chaque métrique, permettant de dire
si l'écart entre deux modèles est significatif. Le coût, à savoir $k$
entraînements au lieu d'un, resterait négligeable ici (le pipeline complet dure
une vingtaine de secondes).

\paragraph{Recherche d'hyperparamètres élargie.}
Les grilles utilisées sont volontairement compactes (quatre valeurs par
hyperparamètre). Une grille plus fine (profondeurs intermédiaires,
\texttt{min\_samples\_leaf} variable, $k$ pairs/impairs, seuils au pas de
$0{,}05$) pourrait apporter quelques points supplémentaires, au prix d'un
temps de calcul encore raisonnable.

\paragraph{Forêt aléatoire.}
L'extension la plus naturelle du travail : entraîner plusieurs arbres CART sur
des échantillons bootstrap avec sous-échantillonnage des variables, et moyenner
leurs prédictions. Notre implémentation de CART constitue déjà la brique de
base ; la forêt réduirait l'instabilité de l'arbre seul, sa principale
faiblesse structurelle.

\paragraph{Traitement explicite du déséquilibre.}
Compléter le réglage du seuil par une pondération des classes dans la fonction
de coût, ou par un ré-échantillonnage de l'entraînement (sur-échantillonnage
des défauts ou sous-échantillonnage des remboursements), pourrait relever le
rappel (la métrique critique) au prix d'un arbitrage à documenter sur la
précision.

\paragraph{Choix du seuil par coûts métier.}
Le F1-score accorde le même poids à la précision et au rappel. Une banque
préférerait minimiser un coût total explicite, de la forme
$C = c_{FN} \cdot FN + c_{FP} \cdot FP$ avec $c_{FN} \gg c_{FP}$ : le seuil
optimal s'en déduirait directement, et la comparaison des modèles se ferait en
euros plutôt qu'en points de F1.

\paragraph{Visualisations complémentaires.}
Deux figures enrichiraient l'analyse : la courbe précision-rappel, plus
informative que la ROC en présence de déséquilibre, et l'importance des
variables de l'arbre (réduction totale de Gini par variable), qui relierait
les performances aux intuitions métier de la section~\ref{sec:donnees}.
